\documentclass[11pt]{article}

% ===================== PACKAGES =====================
\usepackage[a4paper,margin=1in]{geometry}
\usepackage{amsmath,amssymb}
\usepackage{enumitem}
\usepackage{fancyhdr}
\usepackage{xcolor}

% ===================== HEADER & FOOTER =====================
\pagestyle{fancy}
\fancyhf{}
\lhead{CSE 400: Fundamentals of Probability in Computing}
\rhead{Lecture 21 Scribe}
\cfoot{\thepage}

% ===================== TITLE =====================
\title{
    \normalsize School of Engineering and Applied Science (SEAS), Ahmedabad University \\
    \vspace{0.2cm}
    \textbf{CSE 400: Fundamentals of Probability in Computing}\\
    \Large Lecture 21 Scribe 
}
\author{}
\date{}

\begin{document}
\maketitle

\vspace{-2cm}
\begin{center}
    \begin{tabular}{ll}
        \textbf{Name:} & {Khushi Paghdar \hspace{2.5in}} \\ [1.5ex]
        \textbf{Enrollment No:} & {AU2440131 \hspace{2.5in}} \\ [1.5ex]
        \textbf{Email:} & {khushi.p12@ahduni.edu.in \hspace{2.1in}} \\ [1.5ex]
        \textbf{Lecture:} & {21 \hspace{2.9in}} \\ [1.5ex]
        
    \end{tabular}
\end{center}

\hrule
\vspace{0.5cm}

% ===================== CONTENT =====================

\section*{Objective}
This scribe serves as an academic reference for exam preparation for Lecture 21.

\section{Introduction to Tail Probabilities}

\subsection{What is a Tail?}

Let $X$ be a random variable. A tail probability refers to the probability that $X$ takes values far away from typical or expected values, specifically:
\[
P\{X > x\}
\]

This represents the right tail of the distribution.

\subsection*{Step-by-Step Understanding}
\begin{enumerate}
    \item A random variable $X$ has a distribution over possible values.
    \item Most values cluster around the mean $E[X]$.
    \item The tail region consists of values much larger than the mean.
    \item The expression $P\{X > x\}$ measures how likely it is that $X$ exceeds a threshold $x$.
\end{enumerate}

\subsection*{Interpretation}
\begin{itemize}
    \item If $x$ is large, then $P\{X > x\}$ should ideally be small.
    \item This probability quantifies extreme events.
\end{itemize}

\section{Motivation for Tail Bounds}

\subsection{Why Study Tail Bounds?}

In many systems (especially computing systems), we are interested in bounding the probability of extreme events. Tail bounds help us:
\begin{itemize}
    \item Guarantee system performance
    \item Ensure reliability
    \item Limit the probability of undesirable outcomes
\end{itemize}

\section{Markov’s Inequality}

\subsection{Statement}

Let $X$ be a non-negative random variable. Then for any $a > 0$:
\[
P(X \geq a) \leq \frac{E[X]}{a}
\]

\subsection{Conditions}
\begin{itemize}
    \item $X \geq 0$
    \item $a > 0$
\end{itemize}

\section{Proof of Markov’s Inequality}

\subsection*{Step-by-Step Proof}

\begin{enumerate}
    \item Start with expectation:
    \[
    E[X] = \sum x \cdot P(X = x)
    \]

    \item Split the sum into two parts:
    \[
    E[X] = \sum_{x < a} xP(X=x) + \sum_{x \geq a} xP(X=x)
    \]

    \item Since $x \geq a$ in second sum:
    \[
    \sum_{x \geq a} xP(X=x) \geq \sum_{x \geq a} aP(X=x)
    \]

    \item Therefore:
    \[
    E[X] \geq a \cdot P(X \geq a)
    \]

    \item Divide both sides by $a$:
    \[
    P(X \geq a) \leq \frac{E[X]}{a}
    \]
\end{enumerate}

\section{Key Insights}

\begin{itemize}
    \item Provides an upper bound using only the mean.
    \item Does not require full distribution knowledge.
    \item Works only for non-negative random variables.
\end{itemize}

\section{Conclusion}

Markov’s Inequality is a fundamental tool for bounding tail probabilities. It provides a simple yet powerful way to estimate the likelihood of extreme events using only expected value.

\bigskip
\hrule
\vspace{0.2cm}
\begin{center}
    \small \textit{End of Submission}
\end{center}

\end{document}